\providecommand{\Res}[2]{\qty[round-mode=figures,round-precision=3]{\fpeval{#1}}{#2}}
\providecommand{\Val}[1]{\num[round-mode=figures,round-precision=4]{\fpeval{#1}}}
\providecommand{\ResSci}[2]{\qty[round-mode=figures,round-precision=3,exponent-mode=scientific]{\fpeval{#1}}{#2}}
\subsubsection*{Question a) Vitesse de la particule alpha}
\begin{enumerate}[label=\arabic*)]
    \item \textbf{Variables connues et recherchées :} $K_\alpha = \qty{\v{K}}{MeV} = \ResSci{\vsi{K}}{J}$, $m_\alpha = \qty{6.645e-27}{kg}$ ; variable recherchée : $v_\alpha$
    \item \textbf{Équations pertinentes :}
    \[ K = \frac{1}{2} m v^2 \]
    \item \textbf{Isolation de la variable recherchée :}
    \[ v_\alpha = \sqrt{\frac{2K_\alpha}{m_\alpha}} \]
    \item \textbf{Résolution numérique :}
    \[ v_\alpha = \sqrt{\frac{2 \cdot \ResSci{\vsi{K}}{}}{\num{6.645e-27}}} = \ResSci{sqrt(2*\vsi{K}/6.645e-27)}{m/s} \]
    \item \textbf{Réponse :} $v_\alpha = \mathbf{\ResSci{sqrt(2*\vsi{K}/6.645e-27)}{m/s}}$ (\Res{100*sqrt(2*\vsi{K}/6.645e-27)/3e8}{\percent} de $c$, l'approche classique est acceptable).
\end{enumerate}

\subsubsection*{Question b) Énergie de recul}
\begin{enumerate}[label=\arabic*)]
    \item \textbf{Variables connues et recherchées :} masses dans le rapport $m_f / m_\alpha = (A-4)/4 = \fpeval{\v{A}-4}/4$ ; variable recherchée : $K_f$
    \item \textbf{Équations pertinentes :}
    \[ m_\alpha v_\alpha = m_f v_f \qquad K = \frac{p^2}{2m} \]
    \item \textbf{Isolation de la variable recherchée :} les deux quantités de mouvement ont le même module, donc
    \[ \frac{K_f}{K_\alpha} = \frac{m_\alpha}{m_f} \implies K_f = K_\alpha\,\frac{4}{A - 4} \]
    \item \textbf{Résolution numérique :}
    \[ K_f = \num{\v{K}} \cdot \frac{4}{\fpeval{\v{A}-4}} = \Res{\v{K}*4/(\v{A}-4)}{MeV} \]
    \item \textbf{Réponse :} le noyau fils recule avec \textbf{\Res{1000*\v{K}*4/(\v{A}-4)}{keV}}.
\end{enumerate}

\subsubsection*{Question c) Énergie libérée}
\[ Q = K_\alpha + K_f = \num{\v{K}} + \Val{\v{K}*4/(\v{A}-4)} = \Res{\v{K}*\v{A}/(\v{A}-4)}{MeV} \]
La désintégration libère \textbf{\Res{\v{K}*\v{A}/(\v{A}-4)}{MeV}} ; la particule alpha en emporte presque toute l'énergie.
